% ==============================================================================
% === TEMPLATE DE CONTENIDO - Wasaff Consulting                           ===
% === Este es el esqueleto del informe. Carga el estilo y define las secciones. ===
% ==============================================================================

\documentclass[11pt, a4paper]{article}

% --- CARGA DE LA PLANTILLA DE ESTILO ---
\usepackage{wasaff_consulting}
\usepackage{lipsum} % Paquete para generar texto de relleno (Lorem Ipsum)

% --- METADATOS DEL DOCUMENTO (Se definen en main.tex) ---
% Dejamos estos comandos aquí como recordatorio de lo que se debe definir
% en el archivo principal.
% \title{Título del Proyecto}
% \author{Felipe Wasaff}
% \wctitle{...}
% \wcsubtitle{...}
% \wccovergraphic{...}
% \wccredits{...}
% \date{...}

% --- INICIO DEL DOCUMENTO ---
\begin{document}

% Genera la portada personalizada
\makeWCTitlePage

% Índice
\tableofcontents
\clearpage

% --- SECCIONES ESTÁNDAR DEL INFORME ---

\section*{Resumen Ejecutivo}
\lipsum[1] % Párrafo de Lorem Ipsum

\subsection*{Hallazgos Principales}
\begin{itemize}
    \item \lipsum[2][1-2] % 2 frases del párrafo 2
    \item \lipsum[3][1-2]
    \item \lipsum[4][1-2]
\end{itemize}

\begin{figure}[h!]
    \centering
    % \includegraphics[width=0.9\textwidth]{grafico_principal.png}
    \fcolorbox{black}{WCBgGray}{\begin{minipage}[c][6cm][c]{0.9\textwidth}\centering \textbf{GRÁFICO PRINCIPAL}\end{minipage}}
    \caption{Descripción del gráfico principal que resume el hallazgo más importante del estudio.}
\end{figure}

\clearpage

\section*{Análisis de Profundidad}
\lipsum[5]

\begin{figure}[h!]
    \centering
    \begin{minipage}[t]{0.48\textwidth}
        % \includegraphics[width=\linewidth]{grafico_secundario_1.png}
        \fcolorbox{black}{WCBgGray}{\begin{minipage}[c][5cm][c]{\linewidth}\centering \textbf{GRÁFICO 2}\end{minipage}}
        \caption{Descripción del primer gráfico secundario que profundiza en un aspecto del análisis.}
    \end{minipage}
    \hfill
    \begin{minipage}[t]{0.48\textwidth}
        % \includegraphics[width=\linewidth]{grafico_secundario_2.png}
        \fcolorbox{black}{WCBgGray}{\begin{minipage}[c][5cm][c]{\linewidth}\centering \textbf{GRÁFICO 3}\end{minipage}}
        \caption{Descripción del segundo gráfico secundario que explora otra variable o escenario.}
    \end{minipage}
\end{figure}

\subsection*{Nota Metodológica}
\lipsum[6]

\clearpage

\section*{Hoja de Ruta y Recomendaciones}
\lipsum[7][1-3]

\begin{table}[h!]
    \centering
    \renewcommand{\arraystretch}{1.5}
    \begin{tabularx}{\textwidth}{>{\bfseries}p{3cm} X X}
    \toprule
    \rowcolor{WCGray}
    \textbf{\color{white}Oportunidad} & \textbf{\color{white}Acción Propuesta} & \textbf{\color{white}Justificación y ROI} \\
    \midrule
    \rowcolor{WCLightGray}
    \textbf{Recomendación 1} & \lipsum[8][1-2] & \lipsum[9][1-2] \\
    \textbf{Recomendación 2} & \lipsum[10][1-2] & \lipsum[11][1-2] \\
    \rowcolor{WCLightGray}
    \textbf{Recomendación 3} & \lipsum[12][1-2] & \lipsum[13][1-2] \\
    \bottomrule
    \end{tabularx}
\end{table}

\vfill
\begin{center}
    \fcolorbox{WCBlue}{WCBlue}{%
    \href{https://wasaffconsulting.com/como-empezar}{%
        \begin{minipage}{0.6\textwidth}
            \centering
            \color{white}\raleway\Large\bfseries\strut Inicie su Propio Diagnóstico
        \end{minipage}%
    }}
\end{center}

\end{document}
